% Generated by render_comparison.py from saved evidence; do not hand-edit tables.
\subsection{比較對象與共同條件}
本節整理 2026-11-11 的本機三方 TensorRT 研究。程式基準為 \code{3cbaa518f0dc}，新增比較證據未包含在原基準 commit 中。本次只讀取既有分數與測時，不重跑推論、計分、訓練或調參。

三模型為官方預訓練 YOLO26m、UA2 teacher e4 primary EMA 微調版，以及現行 Full25 mixed student e10 primary EMA HR Mamba。

三方同為 native F01／F02、1280×768、batch 1、FP16 allowed、FP32 I/O、共用 F04 與 G1 tracker；Graph ON，slots=2／2／4，NSA 與量測框 OFF。各模型產生自己的偵測及 detector-conditioned GMC，並非固定同一份 motion 的單因素消融。

\begin{finding}
\textbf{官方對照使用 raw6 one2many head，接本專案共用 F04／G1 tracker。}\par
它與 Ultralytics 預設 one2one/end2end 的輸出契約不同，也不是官方追蹤器排名。既有研究另有原始 one2one 的模型層控制，但未重評其完整追蹤品質，故不混入品質表。
\end{finding}

品質指標為 0–100 尺度，差值為百分點（pp）；IDSW 是次數，越低越好。DetA／AssA 是 HOTA 的偵測／關聯指標，不是 AP50。合併採 frozen TrackEval 的完整 raw objects 與 \code{combine_sequences}，不平均各段 HOTA。表格取三位小數，完整精度與 FP／FN／Frag 留在原始 JSON／CSV。\cite{primary156,trtcompare}

\subsection{五資料集覆蓋與整體結果}
每模型覆蓋 156 段、135,503 幀，合計 468 cells／406,509 幀；21 個 MOT17 cells 逐檔驗證後沿用前輪 native TRT，其餘 447 cells 於該輪推論。PP22 為 PersonPath22，48,452 是完整序列幀數，其中僅 8,753 個有 GT 的原始影格參與計分；其餘幀仍推論，計分前只保留 GT 存在幀的預測與原幀號。這是本機指定 split／資料池，並非官方隱藏測試集排名。

來源為 MOT17 train 七個不同 SDP 序列、MOT20 train 四段、MOT15 train 十一段、DanceTrack train 40＋val 25 段、PP22 mot\_test\_kf 69 段。Train 指原始資料集目錄名，實際本地模型訓練角色另述。
\begin{table}[H]
\caption{各資料集 HOTA 與樣本量；Δ 為 HR 減比較對象（pp）}
\small
\noindent\begin{tabularx}{\textwidth}{@{}Yrrrrrrr@{}}
\toprule
資料集 & 段數 & 幀數 & 官方 & 微調 & HR & Δ 官方 & Δ 微調 \\
\midrule
MOT17 & 7 & 5,316 & 46.011 & 61.030 & 63.083 & +17.072 & +2.053 \\
MOT20 & 4 & 8,931 & 16.892 & 29.647 & 37.390 & +20.498 & +7.742 \\
MOT15 & 11 & 5,500 & 43.277 & 44.901 & 45.054 & +1.777 & +0.153 \\
DanceTrack & 65 & 67,304 & 35.869 & 36.245 & 35.702 & -0.167 & -0.543 \\
PP22 & 69 & 48,452 & 43.638 & 43.623 & 41.649 & -1.989 & -1.974 \\
\bottomrule
\end{tabularx}
\end{table}
\source{primary156 的 \path{metrics/summary.json} 及獨立 scoring audit\cite{primary156}。}
\subsection{七個主要群組與逐段差異}
\begin{table}[H]
\caption{主要群組 HOTA；勝／平／負為逐段 HR 相對微調版}
\small
\noindent\begin{tabularx}{\textwidth}{@{}Yrrrrrr@{}}
\toprule
群組 & 段數 & 官方 & 微調 & HR & Δ 微調 & 勝／平／負 \\
\midrule
MOT17/all & 7 & 46.011 & 61.030 & 63.083 & +2.053 & 4/0/3 \\
MOT20/all & 4 & 16.892 & 29.647 & 37.390 & +7.742 & 4/0/0 \\
MOT15/moving & 5 & 40.244 & 44.184 & 43.335 & -0.849 & 1/0/4 \\
MOT15/other & 6 & 47.116 & 45.783 & 47.275 & +1.492 & 1/0/5 \\
DanceTrack/all & 65 & 35.869 & 36.245 & 35.702 & -0.543 & 29/0/36 \\
PP22/moving & 9 & 45.152 & 44.694 & 42.607 & -2.087 & 2/0/7 \\
PP22/static & 60 & 43.527 & 43.546 & 41.581 & -1.965 & 14/0/46 \\
\bottomrule
\end{tabularx}
\end{table}
HR 相對微調 YOLO26m 的七個主要群組為 3 勝／0 平／4 負，156 段為 55 勝／0 平／101 負。MOT17／MOT20 合併 HOTA 分別增加 2.053／7.742 pp；DanceTrack 與 PP22 則分別下降 0.543／1.974 pp。MOT15 合併僅增加 0.153 pp，moving 子群仍下降 0.849 pp，且全 11 段只有 2 段增加。即使相對官方預訓練版，HR 在 DanceTrack 與 PP22 的 HOTA 也較低。因此現有結果支持資料集相關的取捨，不能宣稱普遍優於官方或微調原版。

MOT17 雖然 HOTA／IDF1 合併增加，IDSW 仍由微調版 364 增至 HR 385；MOT20 由 2,401 增至 3,754，PP22 由 4,484 增至 5,138。需同時檢視偵測、關聯與身份切換，不能從單一總分推論身份穩定全面改善。

\subsection{偵測、關聯與身份指標}
\begin{table}[H]
\caption{三模型完整品質結果；各資料集採正式合併方法}
\small
\noindent\begin{tabularx}{\textwidth}{@{}P{22mm}Yrrrrrr@{}}
\toprule
資料集 & 模型 & HOTA & DetA & AssA & IDF1 & MOTA & IDSW \\
\midrule
MOT17 & 官方 YOLO26m & 46.011 & 42.638 & 50.107 & 55.597 & 45.275 & 332 \\
MOT17 & 微調 YOLO26m & 61.030 & 63.036 & 59.670 & 71.329 & 73.291 & 364 \\
MOT17 & 現行 HR Mamba & 63.083 & 66.295 & 60.523 & 71.528 & 76.216 & 385 \\
MOT20 & 官方 YOLO26m & 16.892 & 8.196 & 34.995 & 15.976 & 9.598 & 418 \\
MOT20 & 微調 YOLO26m & 29.647 & 27.800 & 32.076 & 41.116 & 36.965 & 2401 \\
MOT20 & 現行 HR Mamba & 37.390 & 41.121 & 34.505 & 51.204 & 54.961 & 3754 \\
MOT15 & 官方 YOLO26m & 43.277 & 40.579 & 46.502 & 55.483 & 22.676 & 261 \\
MOT15 & 微調 YOLO26m & 44.901 & 41.253 & 49.103 & 56.522 & 19.777 & 211 \\
MOT15 & 現行 HR Mamba & 45.054 & 41.218 & 49.557 & 56.816 & 19.178 & 211 \\
DanceTrack & 官方 YOLO26m & 35.869 & 60.478 & 21.515 & 36.675 & 74.088 & 7734 \\
DanceTrack & 微調 YOLO26m & 36.245 & 61.008 & 21.775 & 37.520 & 76.594 & 7646 \\
DanceTrack & 現行 HR Mamba & 35.702 & 59.367 & 21.702 & 37.431 & 73.996 & 8318 \\
PP22 & 官方 YOLO26m & 43.638 & 45.211 & 42.521 & 48.216 & 47.942 & 4517 \\
PP22 & 微調 YOLO26m & 43.623 & 45.663 & 42.240 & 50.269 & 43.333 & 4484 \\
PP22 & 現行 HR Mamba & 41.649 & 44.103 & 39.914 & 47.947 & 39.138 & 5138 \\
\bottomrule
\end{tabularx}
\end{table}
\subsection{同條件速度與成本取捨}
此表引用另一輪受控效能研究：RTX 5070 Ti Laptop、相同 TensorRT SDK／輸入與交錯順序；tracking 使用 MOT17-10 前 140 幀，20 幀暖機、每模式四個 fresh sessions／480 個量測樣本，JPEG 預載。Single 為 enqueue\ensuremath{\to}receive；queued FPS 為整段完成吞吐。均排除讀檔、初始化、暖機與寫檔，不代表完整資料集或攝影機端到端效能。F03 為模型層 GPU events，與完整 tracking 的量測終點不同。

\begin{table}[H]
\caption{MOT17-10 prefix 的受控效能；F03 與 tracking 分開量測}
\small
\noindent\begin{tabularx}{\textwidth}{@{}Yrrrrr@{}}
\toprule
模型 & F03 ms & Single mean ms & Single p95 ms & Queued FPS & Engine MiB \\
\midrule
官方 YOLO26m & 3.744 & 5.734 & 6.938 & 245.175 & 42.506 \\
微調 YOLO26m & 3.809 & 5.938 & 7.178 & 230.123 & 40.868 \\
現行 HR Mamba & 4.790 & 6.816 & 7.139 & 194.417 & 57.939 \\
\bottomrule
\end{tabularx}
\end{table}
HR 相對微調版在 MOT17-10 的 single 平均延遲增加 14.79\%、queued 吞吐下降 15.52\%；另測 MOT17-02 為 +17.42\%／−15.13\%。F03 Graph mean 增加 25.76\%，engine 檔案增加 41.77\%；execution-context scratch bound 同為約 135.17 MB，不代表完整行程顯存相同。MOT17-10 的 HR single p95 略低，02 則較高，尾延遲亦非全場景一致。

另列官方原始 one2one/end2end 的模型層 Graph 控制，四張固定輸入、每模型四個 fresh processes／1,280 樣本。
\begin{table}[H]
\caption{官方原始輸出契約的模型層控制；輸出終點不同}
\small
\noindent\begin{tabularx}{\textwidth}{@{}Yrr@{}}
\toprule
模型 & GPU mean ms & GPU p95 ms \\
\midrule
官方原始 one2one & 3.862 & 3.916 \\
現行 HR Mamba & 4.871 & 4.955 \\
\bottomrule
\end{tabularx}
\end{table}
官方輸出為 \code{[1,300,6]}，engine 內含 decode／TopK；HR 仍為 raw6。此控制未重評完整 tracking 品質或吞吐，亦不與前表不同 campaign 的均值相減推導階段成本。
\source{TensorRT 性能比較、model／pipeline summaries 與 model-memory.json\cite{trtcompare}。}
\subsection{資料曝光、數值與驗收邊界}
\begin{itemize}
\item 全部 156 段皆已曝露；MOT17 的 04／05／09／11／13 為兩條本地模型 lineage 的訓練序列、10 用於選模，02 是 M4 之後重用的原 held-out。官方預訓練 corpus 未全面稽核，不能套用本地訓練分工。DanceTrack 為 train＋val 的既定資料池，其餘來源／split 見 frozen selection。
\item 微調版與 HR 的訓練配方、來源及 score calibration 不同；差值不能單獨歸因於 Mamba。七組與逐段結果優先，跨資料集 pooled 值不作主要排名或採用門檻。
\item 官方／微調 FP16 engine 的既有 1\% raw relative-L2 screen 仍為 FAIL，最大 2.406\%／1.537\%；任務分數、finite／repeat／Graph 通過不取代數值驗收，也不升格成正式 session／sanitizer／部署驗收。
\item SportsMOT 依既定政策排除主要推論、評分、gate 與參數選擇；本節的五資料集與七群組均不含它。
\item primary156 不測速，capture wall time 含 lock wait，不用來推算各資料集 FPS。各資料集沒有完整同口徑 AP50 與效能矩陣者不補造數字；現有可比項目以本節為準。部署模型與預設維持不變。
\end{itemize}
\source{原始 selection、training-provenance、evaluation-scope、獨立 scoring audit 與 TRT validation\cite{primary156,trtcompare}。}
